%% =====================================================================
%%  T-21-06  Stirred Water
%%  -------------------------------------------------------------------
%%  One idea per file. Core is mandatory. Slots in canonical order.
%%  Max 700 words. No layout commands. No filler. New source material
%%  for this entry goes in sources/B3/T-21-06.tex -- never by
%%  rewriting this file (docs/RULES.md R0.1).
%%  =====================================================================
%% @kind: topic
%% @id: T-21-06
%% @title: Stirred Water
%% @subtitle: Anger is a strategically counterproductive state --- so arrange it, in the other person
%% @chapter: c21
%% @order: 60
%% @status: stable
%% @sources: B3
%% @updated: 2026-09-19

\topic{T-21-06}{Stirred Water}{Anger is a strategically counterproductive state --- so arrange it, in the other person}

\begin{core}
Anger and emotion are strategically counterproductive: you must stay calm and objective. But if you can make your enemies angry while you stay calm, you gain a decided advantage. Locate the hairline fracture in their vanity, rattle them through it, and the strings swing to your hand --- the agitated opponent reveals, overcommits, and mistakes the trembling for strength.
\end{core}
\begin{mechanism}
B3's law runs on the economics of composure: anger spends the exact assets a contest needs --- judgment, timing, information discipline. The agitated party talks, telegraphs and escalates; the calm one watches, waits and chooses. Hence the offensive doctrine: since your own anger is a liability, invest it in the opponent. The delivery is vanity-targeted --- the taunt crafted from what they cannot laugh off (the tea master's remark that exposed the samurai's reputation as untested), the deliberate slight, the public needling that promises no retaliation. Once stirred, the opponent's mind stops running strategy and starts running grievance; every subsequent move is dictated by the last thing that stung. The defender's half is the harder discipline: the counter is neither suppression (which leaks) nor expression (which spends), but the studied delay --- noticing the hook, naming it to yourself, and answering at the hour of your choosing. The law's reversal notes the whole game is well known: some opponents stir you to make you simple. The sophisticated target is unmoved by both bait and baiter.
\end{mechanism}
\begin{conditions}
Strongest in contests of nerve and attrition --- litigation, rivalry, negotiation, public feuds --- where a single loss of composure can spend months of position. It weakens where anger costs the target nothing (nothing at stake, nothing to lose), in front of audiences that punish the stirrer more than the stirred, and against practitioners of the delay, who convert the bait into a told-you-so.
\end{conditions}
\begin{application}
  \item Know your own crack before anyone else does: the competence question, the loyalty question, the family question. That is where bait will come.
  \item Install the delay: when hooked, name the emotion internally and buy an hour. The hour is where your strategy lives.
  \item As defence-in-public: lower your voice as they raise theirs. The room prices the contrast within a minute.
  \item As offence, spend bait that is true --- the sting that works is the accurate one, delivered without heat and without apology.
  \item After any exchange, log what stirred you. The log is the map of your remaining levers, in other people's hands.
\end{application}
\begin{feedback}
  \item Working (as defence): your adversary's taunts stop landing and start being visible as expenditure. They are paying, you are watching.
  \item Working (as offence): the opponent opens unasked disclosure --- the angry explain themselves at length.
  \item Failing: you stirred them and they walked away from the contest entirely. The bait had no hook in their actual interests.
  \item Failing: you replied to the hook in the same hour. The strings were held; the reply proved it.
\end{feedback}
\begin{failure}
The stirrer's trade has a tail: audiences eventually audit who threw the first stone, and a person known for engineered provocations is never again believed to be provoked. The habit also self-administers --- a practitioner who weaponises vanity daily gets fluent in everyone's cracks, including their own, and lives in a house of marked doors. And against a calm opponent the stirrer's bait is pure exposure: each failed hook teaches the target more about the stirrer than the reverse.
\end{failure}
\begin{countermeasures}
  \item Treat your own flashes of anger as incoming-intelligence alerts: something was aimed. Find the aim before answering the aim.
  \item Watch the professional needler --- someone whose taunts arrive with no exit path for you. Ask what the room is being prepared to sell or excuse (T-17-05).
  \item Practice the hour rule until it is reflex. Most strategic anger is regretted within sixty minutes.
  \item In your own house, audit the culture of provocation: if your team needles each other for performance, you are training the stirred-water reflex against yourselves.
\end{countermeasures}

\sources{B3}
\seesources
\seealso{T-02-01, T-21-05, T-19-03}
